본 논문은 분기 예측기가 없는 Cortex-M0+에서 부호 가변 시프트 연산의 세 구현을 실측
비교하였다. 분기 제거 C 구현은 기준 분기 구현보다 1.2배 느려 ``branchless가 빠르다''는
상식이 역전되었고, 배럴 시프터의 포화 특성을 이용한 수동 어셈블리는 1.5배 빨랐다. 컴파일러
생성 코드 분석을 통해 이 차이가 언어 규칙(미정의 동작), 명령어 집합 구조(2피연산자 파괴형),
레지스터 할당의 세 층위에서 비롯됨을 보였으며, 이러한 컴파일러 패스의 구조적 한계를 LLM
기반 피드백 루프로 예측하고 보완할 수 있음을 본 연구의 수행 과정 자체로 실증하였다. 향후
다양한 커널과 코어로 실험을 확장하고, LLM이 발견한 변환을 컴파일러 패스로 형식화하는
자동화 도구를 연구할 예정이다.
